\subsection{The Swaddling Bands}\label{sec:swaddling}

St Luke tells the birth of the Lord in a single sentence:

\begin{quote}
And she brought forth her first born son and wrapped him up in swaddling clothes
and laid him in a manger: because there was no room for them in the inn.\par
{\smallskip\noindent\small ---Luke 2:7\par}
\end{quote}

\noindent The angel gives the shepherds the same two things as their sign: ``You
shall find the infant wrapped in swaddling clothes and laid in a manger.'' Near
the end of his Gospel Luke tells of another body taken in loving hands, after
the Cross:

\begin{quote}
And taking him down, he wrapped him in fine linen and laid him in a sepulchre
that was hewed in stone, wherein never yet any man had been laid.\par
{\smallskip\noindent\small ---Luke 23:53\par}
\end{quote}

\noindent In the Douay he is wrapped and laid at both ends of his life, for St
Jerome's Latin, though not Luke's Greek, uses one verb for both,
\latin{pannis eum involvit} at Bethlehem and \latin{involvit sindone} at
Calvary. The Fathers did not need the echo to see the Cross in the crib. St Cyril of Jerusalem, in the lecture the
Church reads on the First Sunday, had already set the swaddling clothes of the
first coming against the light that clothes the Lord at the second. St Gregory
Nazianzen set them against the grave:

\begin{quote}
He was wrapped in swaddling clothes---but He took off the swathing bands of the
grave by His rising again. He was laid in a manger---but He was glorified by
Angels, and proclaimed by a star, and worshipped by the Magi.\par
{\smallskip\noindent\small ---St Gregory Nazianzen, Oration 29\par}
\end{quote}

Pope St Gregory the Great, preaching on the Ascension, read a verse of the Song
of Songs that the Church hears again on 21 December: ``The voice of my beloved,
behold he cometh leaping upon the mountains, skipping over the hills.'' Christ,
Gregory said, came to redeem us by leaps, and he named them:
\latin{De cœlo venit in uterum, de utero venit in præsepe, de præsepe venit in
crucem, de cruce venit in sepulcrum}, and last, \latin{de sepulcro rediit in
cœlum}. From heaven he came into the womb, from the womb into the manger, from
the manger to the cross, from the cross into the tomb, and from the tomb he went
back to heaven. Gregory adds that the
Truth made manifest in the flesh took these leaps for our sake, so that we
might run after him, and he cites the bride of the Song:
\latin{Trahe nos post te, curremus in odorem unguentorum tuorum}, draw us after
thee, and we will run in the odour of thy ointments.

St Bede, as the Catena Aurea gives him, saw in the bands themselves what the
Child gives:

\begin{quote}
He who clothes the whole world with its varied beauty, is wrapped up in common
linen, that we might be able to receive the best robe; He by Whom all things are
made, is folded both hands and feet, that our hands might be raised up for every
good work, and our feet directed in the way of peace.\par
{\smallskip\noindent\small ---St Bede, in the Catena Aurea on Luke 2:7\par}
\end{quote}

\noindent The best robe is the one the father in the parable brings out for the
son who comes home. The way of peace is the last line of Zachary's song, sung
on the morning of 24 December.
